% ---------------------------------------------------------------------
%  preamble.tex – cấu hình chung (XeLaTeX)
% ---------------------------------------------------------------------
\usepackage{fontspec}
\setmainfont{Times New Roman}
\setsansfont{Arial}
\setmonofont{Consolas}[Scale=0.9]
\usepackage[vietnamese]{babel}

\usepackage[a4paper,left=3cm,right=2cm,top=2.5cm,bottom=2.5cm]{geometry}
\usepackage{setspace}
\setstretch{1.3}
\setlength{\parindent}{1cm}
\setlength{\parskip}{4pt}

\usepackage{amsmath,amssymb,amsthm,mathtools}
\usepackage{graphicx}
\graphicspath{{figures/}}
\usepackage{booktabs,tabularx,multirow,array,makecell}
\usepackage{caption,subcaption}
\captionsetup{font=small,labelfont=bf,labelsep=period}
\usepackage{float}
\usepackage{enumitem}
\setlist{nosep,leftmargin=*,itemsep=2pt}
\usepackage{xcolor}
\definecolor{husblue}{HTML}{1F4E79}
\definecolor{husred}{HTML}{C00000}
\definecolor{codebg}{HTML}{F5F7FA}
\definecolor{codecomment}{HTML}{2E7D32}
\definecolor{codekw}{HTML}{1F4E79}

% ---- Tiêu đề chương / mục ----
\usepackage{titlesec}
\titleformat{\chapter}[display]
  {\normalfont\bfseries\color{husblue}}
  {\large\MakeUppercase{\chaptertitlename}\ \thechapter}
  {6pt}
  {\LARGE\MakeUppercase}
  [\vspace{2pt}{\color{husblue}\titlerule[1pt]}]
\titlespacing*{\chapter}{0pt}{-20pt}{18pt}
\titleformat{\section}{\normalfont\large\bfseries\color{husblue}}{\thesection}{0.6em}{}
\titleformat{\subsection}{\normalfont\normalsize\bfseries}{\thesubsection}{0.6em}{}
\setcounter{secnumdepth}{3}
\setcounter{tocdepth}{2}

% ---- Header / footer ----
\usepackage{fancyhdr}
\pagestyle{fancy}
\fancyhf{}
\fancyhead[L]{\small\itshape Parallel Prefix Sum Using OpenMP}
\fancyhead[R]{\small\itshape\nouppercase{\leftmark}}
\fancyfoot[C]{\thepage}
\renewcommand{\headrulewidth}{0.4pt}
\fancypagestyle{plain}{\fancyhf{}\fancyfoot[C]{\thepage}\renewcommand{\headrulewidth}{0pt}}

% ---- Mã nguồn ----
\usepackage{listings}
\lstdefinestyle{cstyle}{
  language=C,
  basicstyle=\ttfamily\small,
  keywordstyle=\color{codekw}\bfseries,
  commentstyle=\color{codecomment}\itshape,
  stringstyle=\color{husred},
  backgroundcolor=\color{codebg},
  frame=leftline, framerule=2pt, rulecolor=\color{husblue},
  numbers=left, numberstyle=\tiny\color{gray}, numbersep=6pt,
  xleftmargin=1.5em, framexleftmargin=1em,
  showstringspaces=false, breaklines=true, tabsize=4,
  morekeywords={size_t,elem_t,int64_t,uint64_t,pragma,omp,parallel,for,simd,reduction,
                inscan,scan,inclusive,barrier,single,nowait,schedule,static},
  extendedchars=true, inputencoding=utf8,
}
\lstset{style=cstyle}
\renewcommand{\lstlistingname}{Đoạn mã}
\renewcommand{\lstlistlistingname}{Danh mục đoạn mã}

% ---- Mã giả ----
\usepackage[vlined,ruled,linesnumbered]{algorithm2e}
\SetKwFor{ForPar}{for}{in parallel do}{end}
\SetKwInput{KwIn}{Đầu vào}
\SetKwInput{KwOut}{Đầu ra}
\renewcommand{\algorithmcfname}{Thuật toán}
\SetAlgoCaptionSeparator{.}

% ---- Hình vẽ ----
\usepackage{tikz}
\usetikzlibrary{arrows.meta,positioning,calc,fit,backgrounds,decorations.pathreplacing,shapes.misc}
\usepackage{pgfplots}
\pgfplotsset{compat=1.18}

% ---- Định lý ----
\theoremstyle{definition}
\newtheorem{definition}{Định nghĩa}[chapter]
\newtheorem{example}{Ví dụ}[chapter]
\theoremstyle{plain}
\newtheorem{theorem}{Định lý}[chapter]
\newtheorem{lemma}[theorem]{Bổ đề}
\newtheorem{proposition}[theorem]{Mệnh đề}
\renewcommand{\proofname}{Chứng minh}

% ---- Tài liệu tham khảo ----
\usepackage{csquotes}
\usepackage[backend=biber,style=ieee,sorting=none]{biblatex}
\addbibresource{references.bib}

% ---- Liên kết ----
\usepackage[hidelinks]{hyperref}
\usepackage[english,nameinlink]{cleveref}   % tên tiếng Việt được khai báo tay bên dưới
\crefname{appendix}{Phụ lục}{Phụ lục}
\crefname{enumi}{}{}
\crefname{paragraph}{Mục}{Mục}
\crefname{figure}{Hình}{Hình}
\crefname{table}{Bảng}{Bảng}
\crefname{chapter}{Chương}{Chương}
\crefname{section}{Mục}{Mục}
\crefname{subsection}{Mục}{Mục}
\crefname{equation}{}{}
\crefname{algocf}{Thuật toán}{Thuật toán}
\crefname{lstlisting}{Đoạn mã}{Đoạn mã}
\crefname{theorem}{Định lý}{Định lý}
\crefname{definition}{Định nghĩa}{Định nghĩa}

% ---- Lệnh tiện ích ----
\newcommand{\code}[1]{\texttt{#1}}
\newcommand{\omp}[1]{\texttt{\#pragma omp #1}}
\newcommand{\Oh}{\mathcal{O}}
\newcommand{\Work}{\mathrm{W}}
\newcommand{\Span}{\mathrm{D}}
\newcommand{\figwidth}{0.9\textwidth}
